\documentclass[11pt,a4paper]{article}

\usepackage[T1]{fontenc}
\usepackage[utf8]{inputenc}
\usepackage{amsmath,amssymb,amsthm}
\usepackage{geometry}
\geometry{margin=2.3cm}
\usepackage{booktabs}
\usepackage{graphicx}
\usepackage{hyperref}
\frenchspacing

\newcommand{\statut}[1]{{\footnotesize\textsf{[#1]}}}
\newcommand{\Mpl}{M_{\mathrm{Pl}}}
\newcommand{\Heff}{\hat{H}_{\mathrm{eff}}}
\newcommand{\Hsys}{\hat{H}_{\mathrm{sys}}}

\theoremstyle{definition}
\newtheorem{postulat}{Postulat}

\renewcommand{\thesection}{\Roman{section}}
\renewcommand{\thesubsection}{\thesection.\Alph{subsection}}

\title{\vspace{-1.0cm}\textbf{Unit\'e 3 --- Dualit\'e d'\'echelle compactifi\'ee :\\
une voie g\'eom\'etrique vers l'unification du quantique et de la gravitation}\\[5pt]
\large Un seul champ d'\'echelle, des trous noirs \`a l'inflation, et un pont vers l'existence quantique intermittente}
\author{Hassane Bakkaoui\\ \small Chercheur ind\'ependant --- bakkahassa@hotmail.com}
\date{\small Juin 2026 $\cdot$ synth\`ese (gr-qc / astro-ph.CO / quant-ph)}

\begin{document}
\maketitle

\begin{abstract}
\noindent
L'unification de la m\'ecanique quantique et de la relativit\'e g\'en\'erale reste ouverte.
Nous proposons que leurs deux faces \'emergent d'une m\^eme structure : la compactification
des axes d'\'echelle, spatial et temporel, en cercles ---
$\mathcal{M}=M_{3+1}\times S^1_\psi\times S^1_\chi$ sous $\mathbb{Z}_2\times\mathbb{Z}_2$. Un
seul champ d'\'echelle, coupl\'e non minimalement via $f=1+\beta(\cos\psi+\cos\chi)$ et portant
un potentiel born\'e $V=\alpha(2-\cos\psi-\cos\chi)$, gouverne les deux r\'egimes. Nous
\'etablissons : (i) une inflation naturelle \`a couplage verrouill\'e, pr\'edictive ---
$(n_s,r)$ et une contrainte sur $\beta$ ; (ii) l'absence de tout trou noir r\'egulier dans le
secteur \emph{statique} (cl\^oture), robuste au d\'egel de l'axe temporel ; (iii) en
effondrement, que le verdict --- la courbure est-elle plafonn\'ee ? --- est
\emph{d\'ependant de la r\'eduction} et ouvert : il tient \`a la nature du p\^ole $\infty^*$
(\S\ref{sec:pole}) ; (iv) un pont reliant l'axe temporel $\chi$ au temps quantique de
l'hypoth\`ese de discontinuit\'e d'existence quantique (HDEQ), la gravit\'e modulant le temps.
La m\'etrique d'espace de champ est d\'efinie positive ($f>0$, B1), la bande viable l'intervalle
ouvert $\beta\in(-0{,}5;-0{,}15)$ ; les axes d'\'echelle sont des cercles internes
($\sigma$-mod\`ele), sans dimension de genre temps ni fant\^ome (F2) ; la propagation libre ne
viole pas Lorentz. Le facteur d'existence est d\'eriv\'e exactement, $\varepsilon=E(k)/K(k)$.
$\alpha$ est fix\'e par l'amplitude scalaire $A_s$ ($\alpha\approx0{,}5$--$1{,}2\times10^{-9}\,\Mpl^4$,
\'echelle GUT), d'o\`u $\tau_0\approx3{,}5\times10^{-37}$~s \emph{au $f_s\approx5$--$6$ de
l'inflation} et un running $dn_s/d\ln k\approx-6{,}5\times10^{-4}$ (A1) ; le r\'e\'echelonnement
$\Heff=\varepsilon\Hsys$ est montr\'e via une horloge relationnelle de Page--Wootters (D1) ; la
modulation gravitationnelle $\tau_0(R)$ se r\'eduit en laboratoire \`a la dilatation
temporelle de la RG, mais $\tau_0\to\infty$ \`a la courbure d'effondrement $R_c$ (D4).
\textbf{Honn\^etet\'e sur le p\'erim\`etre.} Le secteur inflationnaire reproduit, sans
l'am\'eliorer, l'inflation naturelle non minimale p\'eriodique de la litt\'erature ; la
nouveaut\'e n'est pas l\`a. Le test rigoureux (E1) montre le secteur horloge non discriminant
($\varepsilon=1-10^{-112}$) : le CMB est le seul juge, et la signature propre du cadre
\emph{serait} la corr\'elation inter-secteurs $\beta(\mathrm{CMB})=\beta(\text{cl\^oture})$ ---
une conjecture conditionnelle \`a la r\'eduction (\S\ref{sec:pole}), non un th\'eor\`eme. Chaque
\'enonc\'e est \'etiquet\'e selon son statut.

\smallskip\noindent\textbf{Mots-cl\'es :} dualit\'e d'\'echelle ; gravit\'e scalaire-tenseur ;
couplage non minimal ; inflation naturelle ; trous noirs r\'eguliers ; existence quantique
intermittente ; d\'ecoh\'erence intrins\`eque ; horloge relationnelle ; unification.\\
\textbf{PACS :} 04.50.Kd, 04.70.Bw, 98.80.Cq, 03.65.Ta, 04.62.+v.\quad
\textbf{MSC :} 83C57, 83D05, 83F05, 81P15.

\smallskip\noindent\textit{Statut \'epist\'emique :} \statut{preuve} alg\`ebre exacte / calcul
formel $\cdot$ \statut{num\'erique} int\'egration reproductible $\cdot$ \statut{conjecture}
plausible, non prouv\'e $\cdot$ \statut{heuristique} analogie $\cdot$ \statut{sp\'eculation}
ouvert $\cdot$ \statut{ouvert} \`a \'etablir $\cdot$ \statut{\`a v\'erifier} r\'ef.\ externe.
Aucune conjecture n'est pr\'esent\'ee comme un th\'eor\`eme.
\end{abstract}

% =====================================================================
\section{Introduction}
R\'econcilier m\'ecanique quantique et relativit\'e g\'en\'erale est un probl\`eme ouvert depuis
un si\`ecle. C\^ot\'e gravitation, deux exigences p\`esent sur toute physique du r\'egime de
Planck : faire dispara\^itre les singularit\'es [3--7] et rendre compte de l'expansion
acc\'el\'er\'ee primordiale [8, 9]. C\^ot\'e quantique, la d\'ecoh\'erence et la s\'election d'une
base privil\'egi\'ee doivent \^etre expliqu\'ees sans environnement externe ad hoc [2, 20]. Les
programmes \'etablis (cordes, gravit\'e \`a boucles, g\'eom\'etrie non commutative, gravit\'e
semi-classique hybride) abordent ces fronts par des voies distinctes.

Nous explorons une unique hypoth\`ese g\'eom\'etrique. L'axe d'\'echelle logarithmique,
habituellement non compact, est pris pour un cercle : par projection st\'er\'eographique,
l'\'echelle de Planck (p\^ole $P$) et une \'echelle « unifi\'ee » (p\^ole $\infty^*$) deviennent
deux p\^oles d'une m\^eme vari\'et\'e [1]. Il y a deux tels axes --- spatial ($\psi$) et temporel
($\chi$) --- d'o\`u l'ar\`ene $\mathcal{M}=M_{3+1}\times S^1_\psi\times S^1_\chi$. L'axe spatial
gouverne la face gravitationnelle/cosmologique [1] ; l'axe temporel, la face quantique --- le
temps de l'hypoth\`ese de discontinuit\'e d'existence quantique (HDEQ) [2].

\paragraph{Une clarification ontologique, d'embl\'ee.} L'HDEQ ne postule pas un temps discret.
Le temps reste un param\`etre continu ; ce qui est intermittent, c'est l'existence active des
syst\`emes quantiques, en cycles de dur\'ee $\tau_0$ (fraction active $\varepsilon$). Le
qualificatif « temporel » pour l'axe $\chi$ d\'esigne le r\^ole d'horloge qu'il joue
(\S\ref{sec:quantique}), non une signature m\'etrique de genre temps.

\paragraph{P\'erim\`etre et honn\^etet\'e \'epist\'emique.} Ce travail ne pr\'etend pas r\'esoudre
l'unification \`a lui seul. Il propose une direction coh\'erente et falsifiable : les choix de
mod\'elisation sont assum\'es, les r\'esultats \'etiquet\'es selon leur statut, et c'est \`a la
communaut\'e de tester, \'etendre ou r\'efuter.

% =====================================================================
\section{Cadre formel : l'ar\`ene et l'action}
\begin{postulat}[compactification]
La vari\'et\'e fondamentale est $\mathcal{M}=M_{3+1}\times S^1_\psi\times S^1_\chi$, les axes
d'\'echelle \'etant des cercles via $\psi=2\arctan[\ln(\lambda/\ell_P)]$ et
$\chi=2\arctan[\ln(\tau/\tau_P)]$. Les p\^oles $\pm\pi$ sont identifi\'es ($\infty^*$) ; $0$ est
le point de Planck $P$.
\end{postulat}

\subsection{Action, m\'etrique d'espace de champ (B1) et bande viable}
\begin{equation}
S=\int d^4x\,\sqrt{-g}\Big[\tfrac12 fR-\tfrac12\omega\big((\partial\psi)^2+(\partial\chi)^2\big)-V\Big],
\quad f=1+\beta(\cos\psi+\cos\chi),\quad V=\alpha(2-\cos\psi-\cos\chi).
\end{equation}
\statut{preuve} (B1) Au repère d'Einstein, la m\'etrique d'espace de champ $\mathcal{G}_{ab}$
est d\'efinie positive $\Leftrightarrow f>0$ : une partie cin\'etique $\omega/f$ (d\'efinie
positive) plus une partie non minimale de rang 1 semi-d\'efinie positive. La bande viable est
l'intervalle ouvert $\beta\in(-0{,}5;-0{,}15)$ (\`a $\beta=-1/2$, $f\to0$ en $P$).
\statut{preuve} (F2) $S^1_\psi,S^1_\chi$ sont des cercles internes (cible $\sigma$-mod\`ele
euclidienne), non des dimensions d'espace-temps ; $\chi$ n'est pas une coordonn\'ee de genre
temps $\to$ pas de fant\^ome, pas de courbe temporelle ferm\'ee. L'enroulement $=$ la charge
d'\'echelle topologique.

% =====================================================================
\section{Le secteur gravitationnel et cosmologique (« l'infiniment grand »)}\label{sec:cosmo}
\statut{preuve} Aux champs gel\'es, la gravit\'e d'Einstein \'emerge ; au vide sym\'etrique $P$
($V=0$), la constante cosmologique s'annule. Au p\^ole $\infty^*$, la source se r\'eduit \`a une
\'energie de vide pure, donnant un espace localement de Sitter de taux
$H^2=2\alpha/[3(1-\beta)]$ (un axe) ou $4\alpha/[3(1-2\beta)]$ (deux axes), avec scalaire de
Kretschmann fini $K=24H^4$.

En cosmologie homog\`ene, le champ canonique $\varphi=f_s\psi$ voit le potentiel
$V=\alpha[1-\cos(\varphi/f_s)]$ : une inflation naturelle \`a couplage non minimal verrouill\'e
--- $f=1+\beta\cos\psi$ a la m\^eme structure que le potentiel (\'etudi\'e ind\'ependamment dans
[16]). \statut{num\'erique} L'int\'egration part de $\infty^*$ (phase de Sitter) et roule
jusqu'\`a $P$ ($\Lambda=0$), avec sortie gracieuse. Pr\'edictions de roulement lent (repère
d'Einstein) :

\begin{table}[h]\centering\small
\begin{tabular}{ccccl}
\toprule
$\beta$ & $f_s$ ($\Mpl$) & $n_s$ & $r$ & statut\\
\midrule
$0$ & $7$ & $0{,}960$ & $0{,}078$ & $r$ trop grand\\
$+0{,}3$ & $7$ & $0{,}949$ & $0{,}128$ & exclu\\
$-0{,}3$ & $5$ & $0{,}958$ & $0{,}017$ & compatible (bord)\\
$-0{,}3$ & $6$ & $0{,}959$ & $0{,}027$ & compatible (bord)\\
$-0{,}3$ & $7$ & $0{,}960$ & $0{,}036$ & au bord BK18\\
\bottomrule
\end{tabular}
\caption{\small Roulement lent $(n_s,r)$ \`a $N=55$. Cible Planck$+$BK :
$n_s=0{,}9649\pm0{,}004$, $r<0{,}036$.}\label{tab:nsr}
\end{table}

Pour le signe naturel $\beta\ge0$, $r\gtrsim0{,}08$ est exclu [12, 13]. La fen\^etre viable est
une bande \'etroite $\beta\in(-0{,}5;-0{,}15)$ avec \textbf{$f_s\approx5$--$6\,\Mpl$}
(intervalle ouvert, cf.\ II.A), donnant $r\approx0{,}02$--$0{,}04$ corr\'el\'e \`a
$n_s\approx0{,}96$. \statut{num\'erique} \emph{Bemol honn\^ete (P1).} Sur cette bande
$n_s\approx0{,}958$--$0{,}961$, soit $\sim1$--$2\sigma$ sous la valeur centrale de Planck : le
cadre vit \emph{au bord} du contour. Les donn\'ees CMB fournissent une contrainte sur le
couplage de dualit\'e d'\'echelle $\beta$ dans ce cadre (la ph\'enom\'enologie $(n_s,r)$ \'etant
partag\'ee avec [16, 24, 25]), testable par Simons Observatory [27] et LiteBIRD [14]
($\sigma(r)\sim10^{-3}$) d'ici 2030.\footnote{Le projet CMB-S4 (DOE--NSF), longtemps
l'instrument de r\'ef\'erence pour ce test, a \'et\'e arr\^et\'e : le 9 juillet 2025, le DOE et la
NSF ont d\'eclar\'e ne plus pouvoir le soutenir (arr\^et ordonn\'e), il n'est donc plus financ\'e
[15]. La falsifiabilit\'e demeure, port\'ee par Simons Observatory (Chili, en service), LiteBIRD
($\sim$2032) et BICEP/Keck (en cours). \statut{\`a v\'erifier}}

\statut{num\'erique} (A1) Le roulement lent au second ordre avec la m\'etrique d'espace de champ
exacte reproduit la Table~\ref{tab:nsr} \`a mieux que $0{,}5\%$, et fixe $\alpha$ par
$A_s=2{,}1\times10^{-9}$ : $\alpha\approx0{,}5$--$1{,}2\times10^{-9}\,\Mpl^4$, soit une \'echelle
d'inflation $V^{1/4}\sim1{,}1$--$1{,}4\times10^{16}$~GeV (GUT). Au $f_s\approx5$--$6$ de
l'inflation, cela donne $\tau_0=2\pi\sqrt{\omega/\alpha}\approx3{,}4$--$3{,}8\times10^{-37}$~s
(coh\'erent : l'horloge utilise d\'esormais le \emph{m\^eme} $f_s$ que l'inflation, ce qui
r\'esout le d\'ecalage ant\'erieur $f_s=57$ donnant $r\approx0{,}14$, exclu). Pr\'ediction stable
sur la bande : le running $dn_s/d\ln k\approx-6{,}5\times10^{-4}$ --- compatible avec Planck
($-0{,}004\pm0{,}007$). (Reliquat : Mukhanov--Sasaki mode-par-mode, sous-pourcent ;
$(n_s,r,\alpha)$ calcul\'es sur la trajectoire repr\'esentative \`a un champ, le traitement
complet \`a deux champs --- offset de $f$ en $\chi=0$, effets de virage --- restant une
hypoth\`ese non test\'ee, piste A2.)

\paragraph{Honn\^etet\'e sur la nouveaut\'e.} Ce secteur reproduit, sans l'am\'eliorer,
l'inflation naturelle non minimale p\'eriodique de Ferreira--Notari--Simeon [16], son extension
scalaire-tenseur [24], et les analyses $(n_s,r,\text{running})$ face \`a BICEP/Keck [25]. La
nouveaut\'e du pr\'esent travail n'est donc pas la ph\'enom\'enologie inflationnaire, mais la
tentative d'unification inter-secteurs (\S\ref{sec:pole}).

% =====================================================================
\section{Le secteur statique : aucun trou noir r\'egulier (cl\^oture)}\label{sec:static}
\statut{preuve}\statut{num\'erique} \`A un champ [1], le profil radial statique est clos : la
branche gel\'ee $\psi\equiv-\pi$ est exactement Schwarzschild--de Sitter ; la branche roulante
d\'eveloppe un horizon singulier (effondrement de $B=e^{\Phi+\Lambda}$, masse de Misner--Sharp
n\'egative, $R$ divergent). Nous appuyons la cl\^oture sur l'argument \emph{robuste} plut\^ot
que sur un tableau hors-bande : par Derrick $+$ topologie, le profil statique (reliant
$\infty^*$ \`a $P$) est un \emph{demi-soliton non prot\'eg\'e} $N=\tfrac12$, donc contractible
--- aucune solution r\'eguli\`ere interpolante n'existe, pour tout $\beta$ de la bande. Le
soliton entier $N=1$ a une masse $8\sqrt{\omega\alpha}\approx2$--$5\times10^{15}$~GeV (\'echelle
GUT ; $\approx3\times10^{15}$ au centre $f_s\approx6$, avec $\alpha$ fix\'e en A1) --- valeur
harmonis\'ee avec le $f_s$ corrig\'e (P1).

\paragraph{D\'egel de l'axe temporel.} \statut{preuve} En rendant $\chi$ dynamique,
$\chi\equiv0$ reste une solution exacte et reproduit la cl\^oture, mais le couplage de $\chi$
\`a $R$ renormalise la gravit\'e : sur la tranche $\chi=0$,
$f_{\mathrm{eff}}=(1+\beta)+\beta\cos\psi$, et le seuil de stabilit\'e se d\'eplace
($\beta_c=-1/4$ pour deux axes). La conclusion de cl\^oture est robuste ; la valeur pr\'ecise du
seuil est d\'ependante de la r\'eduction ($-1/3$ mono-axe, $-1/4$ sur la tranche $\chi=0$, $-1/6$
sur la diagonale --- voir \S\ref{sec:pole}).

% =====================================================================
\section{Le p\^ole $\infty^*$ : une tension structurelle entre cl\^oture et d\'eroulement}\label{sec:pole}
\statut{num\'erique} L'effondrement \emph{dynamique} (non statique) plafonne-t-il la courbure ?
La r\'eponse tient \`a la nature du p\^ole $\infty^*$ --- qui n'est \emph{pas} ind\'ependante de
la r\'eduction. C'est le point central, que nous consignons comme le r\'esultat structurel de
cette synth\`ese : une tension ouverte, non un verdict tranch\'e.

\paragraph{Lecture mono-axe / Jordan.} Sur l'unique axe actif ($\chi$ gel\'e), pour le
potentiel de Jordan $V$, $\infty^*$ est un maximum ($V''(-\pi)=-\alpha<0$). L'effondrement
sph\'erique homog\`ene avec champ roulant pousse alors $\psi$ \emph{\`a travers} $\infty^*$ : la
courbure n'est pas plafonn\'ee, aucun c\oe ur r\'egulier ne se forme. Dans cette lecture, les
deux verdicts --- sortie gracieuse en cosmologie, pas de c\oe ur r\'egulier en effondrement ---
partagent la cause « $\infty^*$ est un sommet ».

\paragraph{Lecture \`a deux champs / repère d'Einstein.} Mais la cosmologie qui fixe $(n_s,r)$
est, stricto sensu, un syst\`eme \`a deux champs sur la diagonale $\psi=\chi$.
\statut{num\'erique} Dans le paysage diagonal en repère d'Einstein $U=V/f^2$, la nature de
$\infty^*$ change de signe \`a $\beta=-1/6$ : pour $\beta>-1/6$ c'est un maximum (roule), mais
pour $\beta<-1/6$ --- soit \emph{toute la bande viable} $\beta\in(-0{,}5;-0{,}15)$ --- $\infty^*$
est un \emph{minimum local pi\'eg\'e (un faux vide)} derri\`ere une barri\`ere interne
$\theta_b$. Dans cette lecture, un effondrement poussant le champ vers $\infty^*$ pourrait l'y
\emph{retenir} (c\oe ur de Sitter) plut\^ot que de le traverser --- ce qui pointerait au
contraire vers une courbure plafonn\'ee, soit un c\oe ur r\'egulier \emph{possible}.

\begin{figure}[h]\centering
\includegraphics[width=0.74\textwidth]{figs_u3/fig_pole.png}
\caption{\small \statut{num\'erique} Paysage diagonal en repère d'Einstein
$U(\theta)/\alpha=2(1-\cos\theta)/(1+2\beta\cos\theta)^2$ ($\psi=\chi=\theta$, forme close
exacte). Pour $\beta<-1/6$ (vert, bleu --- toute la bande viable) $\infty^*$ ($\theta=\pi$) est
un \emph{minimum pi\'eg\'e} derri\`ere une barri\`ere (faux vide) ; pour $\beta\ge-1/6$ (gris,
rouge) c'est un maximum (le champ traverse). La cl\^oture veut $\infty^*$ pi\'eg\'e, le
d\'eroulement le veut traversable : le m\^eme p\^ole ne peut \^etre les deux.}\label{fig:pole}
\end{figure}

\paragraph{La tension.} Les deux lectures divergent sur la nature m\^eme de $\infty^*$, donc
sur la r\'egularisation de l'effondrement. Elles ne sont pas cosm\'etiques. La cl\^oture
\emph{statique} repose sur $\infty^*$ pi\'eg\'e (Derrick : le champ ne peut connecter $\infty^*$
\`a $P$, la seule solution statique est le Schwarzschild--de Sitter gel\'e), tandis que l'image
litt\'erale du « d\'eroulement » (le parchemin, le champ s'enroulant \emph{\`a travers}
$\infty^*$) exige $\infty^*$ traversable. Le m\^eme p\^ole ne peut \^etre les deux, et la bande
$\beta<-1/6$ qui donne la cl\^oture est exactement celle qui obstrue le d\'eroulement dynamique.
Nous consignons donc :
\begin{itemize}
\item la cl\^oture \textbf{statique} (aucun trou noir r\'egulier) est robuste --- elle d\'ecoule
de Derrick $+$ topologie (demi-soliton contractible $N=\tfrac12$) et tient \`a travers les
r\'eductions. \statut{preuve}
\item le verdict de l'effondrement \textbf{dynamique} est d\'ependant de la r\'eduction et
\emph{ouvert} : traverse-et-non-plafonn\'e dans la lecture mono-axe/$V$, mais faux vide pi\'eg\'e
(donc candidat \`a une courbure plafonn\'ee, un c\oe ur r\'egulier) dans la lecture \`a deux
champs en repère d'Einstein, sur toute la bande viable. \statut{conjecture}\statut{ouvert}
\end{itemize}

\paragraph{Un raccourci qui \'echoue.} On pourrait esp\'erer qu'une vraie dynamique multichamp
(une « spirale » : le champ s'\'ecartant de la diagonale en tournant) supprime $r$ ou soutienne
le d\'eroulement. \statut{num\'erique} Une int\'egration compl\`ete \`a deux champs avec \'elan
lat\'eral inject\'e montre qu'il n'en est rien : le virage est transitoire ($\omega/H\approx1{,}6$
au d\'epart) puis amorti par le frottement de Hubble bien avant la fen\^etre CMB ; $r$ reste
$\approx0{,}046$, jamais supprim\'e. La raison est structurelle --- $P$ est un creux ponctuel,
pas une vall\'ee annulaire, rien ne soutient une orbite. Le « d\'eroulement » se lit donc au
mieux comme un \'enonc\'e \emph{topologique/spatial} (la charge enti\`ere de soliton $N$), et le
d\'eroulement \emph{temporel} litt\'eral (une trajectoire enroul\'ee $N\ge1$ \`a travers
$\infty^*$) comme une signature \emph{conditionnelle} (oscillations log-p\'eriodiques)
disponible seulement \emph{hors} de la bande viable --- non comme un m\'ecanisme dynamique
\'etabli.

\paragraph{La seule ouverture, c\^ot\'e quantique.} S'il existe une r\'egularisation dans la
bande viable, le candidat n'est pas le roulement classique mais l'horloge quantique : \`a la
courbure d'effondrement $R_c=2\alpha/|\beta|$, la masse effective de $\chi$ s'annule et
$\tau_0\to\infty$ (D4, \S\ref{sec:quantique}) --- le cycle ne se referme plus, ce qui pourrait
interdire la travers\'ee classique et retenir le champ en $\infty^*$. C'est la seule voie par
laquelle la boucle d'unification pourrait se fermer. \statut{sp\'eculation}

% =====================================================================
\section{Le secteur quantique et le pont (« l'infiniment petit »)}\label{sec:quantique}
\subsection{Le facteur d'existence d\'eriv\'e exactement (D1)}
\statut{preuve} L'oscillation de $\chi$ dans $V$ est un pendule (sine-Gordon). Pour une
amplitude $\chi_{\max}$ (module elliptique $k=\sin(\chi_{\max}/2)$), l'identit\'e sur
$\langle\mathrm{sn}^2\rangle$ donne la forme exacte (v\'erifi\'ee en int\'egrant
$\chi''=-\sin\chi$ \`a $10^{-16}$) :
\[
\varepsilon=\langle\cos^2(\chi/2)\rangle=\frac{E(k)}{K(k)},\qquad \tau_0(k)=4\sqrt{\omega/\alpha}\,K(k).
\]
\`A petite amplitude $\varepsilon\to1-\chi_{\max}^2/8$ et $\tau_0\to2\pi\sqrt{\omega/\alpha}$.
En s'approchant de la s\'eparatrice ($\infty^*$), $\varepsilon\to0$ (suspension). $\varepsilon$
n'est donc plus un param\`etre libre : c'est $E(k)/K(k)$, fix\'e par l'amplitude.

\begin{figure}[h]\centering
\includegraphics[width=0.66\textwidth]{figs_u3/fig_epsilon.png}
\caption{\small \statut{num\'erique} Le facteur d'existence $\varepsilon=E(k)/K(k)$,
$k=\sin(\chi_{\max}/2)$ (v\'erifi\'e contre le pendule $\langle\cos^2(\chi/2)\rangle$ \`a
$10^{-6}$). \`A petite amplitude $\varepsilon\to1-\chi_{\max}^2/8$ (MQ continue) ; \`a la
s\'eparatrice $\varepsilon\to0$ (suspension). $\varepsilon$ est d\'eriv\'e, pas libre.}\label{fig:eps}
\end{figure}

\subsection{R\'ealisation op\'eratorielle (D1) et modulation gravitationnelle (D4)}
\statut{conjecture} (r\'ealis\'ee, D1) $\chi$ est une horloge relationnelle de Page--Wootters
[22] : sur $\mathcal{H}_{\mathrm{horloge}}\otimes\mathcal{H}_{\mathrm{sys}}$ avec une contrainte
intemporelle, le syst\`eme \'evolue par rapport \`a $\chi$ --- d\'efinissant « temporel » comme
un r\^ole d'horloge, sans signature de genre temps (r\'eponse \`a F2). Pour $\Hsys$
\emph{constant}, l'\'evolution moduldee par $g(\chi)=\cos^2(\chi/2)$ donne sur une p\'eriode
$\hat U(\tau_0)=\exp(-i\Hsys\varepsilon\tau_0)$, soit $\Heff=\varepsilon\Hsys$ --- le r\'esultat
de Floquet $\Omega_F=\varepsilon\Omega_R$ de [2], ici v\'erifi\'e pour $\Hsys$ constant (qubit
modul\'e, erreur $\sim10^{-11}$) ; le cas non commutant (Floquet complet) est estim\'e
n\'egligeable mais \emph{non} prouv\'e. \statut{ouvert} d\'eriver la contrainte intemporelle de
l'action de $\chi$ ; exactitude de la r\'eduction \`a deux \'etats.

\statut{preuve} (D4) Le couplage de $\chi$ \`a $R$ d\'eplace sa masse effective :
$\tau_0(R)=2\pi\sqrt{\omega/(\alpha+\tfrac12\beta\Mpl^2 R)}$. Deux effets, \`a ne pas confondre.
(A) $\sqrt{-g_{00}}\approx1+\phi/c^2$ : dilatation gravitationnelle standard ($\tau_0$ est un
scalaire de temps propre), observable mais non nouvelle --- c'est la RG. (B) le d\'ecalage
dynamique $\delta\tau_0/\tau_0\approx-(\beta/4)(\Mpl^2 R/\alpha)\propto R$ local : il s'annule
dans le vide et vaut $\sim10^{-81}$ dans la mati\`ere --- inobservable. Mais \`a
$R_c=2\alpha/|\beta|$ (le seuil d'effondrement de \S\ref{sec:pole}), la masse effective de
$\chi$ s'annule : $\tau_0\to\infty$, l'horloge quantique g\`ele. Le m\^eme $R_c$ qui d\'eclenche
la descente fait diverger le cycle --- une r\'ealisation concr\`ete du pi\'egeage en $\infty^*$
\statut{conjecture}. ($\tau_0(R)$ est en petite amplitude ; le comportement exact pr\`es de
$R_c$ requiert le potentiel anharmonique complet.)

\begin{figure}[h]\centering
\includegraphics[width=0.66\textwidth]{figs_u3/fig_tau0R.png}
\caption{\small \statut{num\'erique} Modulation gravitationnelle
$\tau_0(R)/\tau_0(0)=1/\sqrt{1-R/R_c}$, $R_c=2\alpha/|\beta|$. La p\'eriode de l'horloge diverge
quand la courbure approche le seuil d'effondrement $R_c$ : l'horloge quantique g\`ele --- la
seule voie candidate pour que le secteur quantique r\'egularise l'effondrement.}\label{fig:tau0R}
\end{figure}

\subsection{S\'ecurit\'e Lorentz du temps quantique}
\statut{preuve} $\tau_0\sim10^{-37}$~s correspond \`a $\hbar/\tau_0\approx10^{13}$~GeV ; il faut
v\'erifier qu'aucune violation de Lorentz n'en r\'esulte (temps de vol des photons, LHAASO). Une
existence intermittente ne viole Lorentz que par une modification de la relation de dispersion
d\'ependant de l'impulsion. Or (i) l'\'evolution propre de $\hat H$ est exactement lin\'eaire en
phase \`a tous ordres en $\tau_0$ (Proposition 1 de l'HDEQ) ; (ii) l'action de Maxwell est
invariante de Weyl en $D=4$, donc le photon est aveugle au facteur conforme $f(\chi)$. La
dispersion du photon libre reste exactement $\omega=|p|c$ : aucun signal LHAASO. La seule
correction (anomalie de trace QED) est logarithmique, non en puissance de $E/E_{\mathrm{QG}}$
--- aucun terme LIV lin\'eaire. \statut{\`a v\'erifier}

\subsection{Test inter-secteurs (E1) : le CMB est le seul juge}
\statut{num\'erique} Sur l'$\alpha$ fix\'e en A1, $\tau_0\in[2{,}5;5{,}3]\times10^{-37}$~s. Point
d\'ecisif : l'unification retire la libert\'e sur $\varepsilon$. Le pont lie $\varepsilon$ au
d\'eplacement actuel du champ, lequel \'egale l'\'energie noire ; num\'eriquement
$\psi_{\mathrm{aujourd'hui}}\sim10^{-56}$~rad, donc $\varepsilon=1-3\times10^{-112}$.
Cons\'equences : l'effet Yb$^+$ $\Delta\Gamma/\Gamma=1-\varepsilon^2\approx10^{-112}$, la
d\'ecoh\'erence intrins\`eque sous le plancher, la r\'esonance \`a $\Delta E\approx1{,}2\times10^{13}$~GeV
--- tout le secteur horloge de l'HDEQ est inobservable. \emph{Bemol (ouvert).} Le cadre pose
$\Lambda=0$ au vide $P$, alors que l'\'energie noire observ\'ee est non nulle ; l'identit\'e
« r\'esidu actuel $=$ \'energie noire » qui fixe $\varepsilon=1-10^{-112}$ n'est pas d\'eriv\'ee,
et la r\'econciliation avec le $\Lambda$ observ\'e reste ouverte (une lecture en quintessence
r\'esiduelle, $\varepsilon$ l\'eg\`erement sous $1$, porterait $\Lambda>0$ au prix d'un
ajustement fin). La conclusion « horloges non discriminantes » tient dans les deux lectures.
\statut{ouvert}

\paragraph{Verdict E1 (rigoureux).} Le cadre passe le test de coh\'erence, mais la confrontation
n'est pas sym\'etrique : le c\^ot\'e horloge est satisfait trivialement ($\varepsilon\approx1$).
Toute la charge de falsifiabilit\'e repose donc sur le CMB ($n_s\approx0{,}96$,
$r\approx0{,}02$--$0{,}04$, running $\approx-6{,}5\times10^{-4}$). Le « m\^eme $(\alpha,\omega)$
pour le CMB et les horloges » est rempli automatiquement c\^ot\'e horloge --- \`a ne \emph{pas}
pr\'esenter comme un test discriminant.

% =====================================================================
\section{Vue d'ensemble et crit\`ere de falsifiabilit\'e}
La vision est un ansatz, non un r\'esultat : une ar\`ene unique dont les lois connues sont ses
deux r\'egimes limites. (\emph{Infiniment grand :}) champs lents $\to$ relativit\'e g\'en\'erale,
$\Lambda=0$ au vide, inflation. (\emph{Infiniment petit :}) dynamique du cercle $\chi\to$
existence intermittente, d\'ecoh\'erence (HDEQ). La charni\`ere est l'axe $\chi$, qui se couple
\`a la courbure et module le temps quantique via la gravit\'e. Les deux faces sont jointes par
un pont d\'esormais partiellement prouv\'e ($\varepsilon$ d\'eriv\'e, modulation incluse ; la
construction op\'eratorielle compl\`ete reste ouverte).

Le crit\`ere de v\'erit\'e est num\'erique, mais asym\'etrique : le CMB
$(n_s,r,\text{running})$ est le juge ; les horloges atomiques fournissent une contrainte de
coh\'erence automatiquement satisfaite ($\varepsilon\approx1$), non un test discriminant. La
signature \emph{propos\'ee} pour ce cadre --- qu'aucun observable cosmologique isol\'e ne
poss\`ede, $(n_s,r,\text{running})$ \'etant partag\'es avec [16, 24, 25] --- est la corr\'elation
inter-secteurs $\beta(\mathrm{CMB})=\beta(\text{cl\^oture})$ : le m\^eme couplage devrait ajuster
l'inflation et satisfaire les seuils de cl\^oture des trous noirs. \emph{Cette corr\'elation est
une conjecture conditionnelle \`a la r\'eduction} (\S\ref{sec:pole}), non un th\'eor\`eme : avec
trois seuils candidats ($-1/3,-1/4,-1/6$), « $\beta(\text{cl\^oture})$ » n'est pas encore un
nombre unique, de sorte que trancher la r\'eduction (piste A2) est le pr\'ealable pour en faire
un v\'eritable discriminant. C'est l\`a que vivrait la falsifiabilit\'e, testable via
LiteBIRD / Simons Observatory.

% =====================================================================
\section{Statut \'epist\'emique et limitations}
\begin{table}[h]\centering\small
\begin{tabular}{p{0.66\textwidth}l}
\toprule
Contenu & Statut\\
\midrule
Action finie ; limite RG ; inflation naturelle \`a couplage verrouill\'e, $(n_s,r)$ et une
contrainte sur $\beta$ ; cl\^oture \emph{statique} robuste (Derrick $+$ topologie) ;
$\mathcal{G}_{ab}$ d\'efinie positive (B1) ; sant\'e $\sigma$-mod\`ele, pas de fant\^ome/CTC
(F2) ; s\'ecurit\'e Lorentz du photon ; $\varepsilon=E(k)/K(k)$ d\'eriv\'e ; $\alpha$ fix\'e par
$A_s$, running $\approx-6{,}5\times10^{-4}$ (A1) ; $f_s\approx5$--$6$ et $\tau_0$ coh\'erents.
& \statut{preuve}\statut{num\'erique}\\
\midrule
$\Heff=\varepsilon\Hsys$ via horloge de Page--Wootters, restreint \`a $\Hsys$ constant (D1) ;
$\tau_0(R)$ d\'eriv\'e, $\tau_0\to\infty$ \`a $R_c$ (D4) ; E1 : secteur horloge non discriminant,
CMB seul juge. & \statut{conjecture} (r\'ealis\'ee)\\
\midrule
\textbf{Tension ouverte centrale :} la nature de $\infty^*$ (traversable vs pi\'eg\'e) gouverne
le verdict \emph{dynamique} d'effondrement et la signature inter-secteurs ; r\'eduction non
tranch\'ee (piste A2). D\'eriver la contrainte de Page--Wootters de l'action de $\chi$ ;
r\'econcilier $\Lambda=0$ vs \'energie noire observ\'ee ; effondrement inhomog\`ene ; pi\'egeage
quantique en $\infty^*$ (r\'egularisation) ; r\`egle de Born ; s\'election d'une branche ;
Mukhanov--Sasaki mode-par-mode. & \statut{ouvert}\statut{sp\'eculation}\\
\bottomrule
\end{tabular}
\end{table}

\noindent La distinction des statuts n'est pas un ornement mais une part constitutive de la
m\'ethode : elle garantit que ce travail peut \^etre r\'eutilis\'e sans qu'aucun \'enonc\'e
d\'epasse ce qui est r\'eellement \'etabli.

% =====================================================================
\section{Conclusion}
Un unique champ de dualit\'e d'\'echelle, sur une ar\`ene \`a axes compactifi\'es (cercles
internes, non des dimensions d'espace-temps), engendre une cosmologie inflationnaire pr\'edictive
et falsifiable, r\`egle la question du trou noir r\'egulier \emph{statique} par une cause
structurelle (Derrick $+$ topologie), et offre un pont vers le temps quantique de l'HDEQ
d\'esormais partiellement prouv\'e : le facteur d'existence $\varepsilon=E(k)/K(k)$ est
d\'eriv\'e, le m\'ecanisme actif/inactif et la propagation libre s\^ure pour Lorentz \'etablis.
La r\'egularisation \emph{dynamique}, en revanche, n'est pas tranch\'ee : elle tient \`a la
nature du p\^ole $\infty^*$, traversable dans une r\'eduction et faux vide pi\'eg\'e dans une
autre --- une tension que nous consignons comme le probl\`eme ouvert central, le gel de
l'horloge \`a $R_c$ \'etant la seule voie candidate pour fermer la boucle. Ce n'est pas une
cl\^oture de l'unification ; c'est une direction dont la nature m\^eme est d'\^etre r\'efutable.
C'est le prix --- l'\'etiquetage honn\^ete de chaque \'enonc\'e --- auquel une intuition
individuelle peut devenir une direction collective.

\paragraph{Remerciements.} L'auteur remercie les outils de calcul formel et num\'erique
(SymPy, SciPy), qui ont permis la v\'erification syst\'ematique \`a chaque \'etape, et la
communaut\'e gr-qc/hep-th/quant-ph pour les standards de rigueur et de falsifiabilit\'e qui ont
guid\'e ce travail. La formalisation a \'et\'e conduite en dialogue avec des syst\`emes
d'assistance analytique, sous la direction de l'auteur, \`a qui revient l'enti\`ere
responsabilit\'e du contenu scientifique.

% =====================================================================
\section*{Reproductibilit\'e}
Les r\'esultats num\'eriques et les v\'erifications symboliques sont reproduits par deux scripts
autonomes (NumPy/SymPy/SciPy/Matplotlib), aux sections cod\'ees pour \^etre propres \`a ce
document (aucune collision avec d'autres jeux de scripts) :
\begin{itemize}
\item[{\texttt{[S1]}}] facteur d'existence $\varepsilon=E(k)/K(k)$ v\'erifi\'e contre le pendule $\langle\cos^2(\chi/2)\rangle$ ; petite amplitude $\varepsilon\to1-\chi_{\max}^2/8$ ; $\tau_0(k)=4\sqrt{\omega/\alpha}\,K(k)$.
\item[{\texttt{[S2]}}] nature de $\infty^*$ : seuil diagonal $\beta=-1/6$ (pi\'eg\'e/roule) et masse transverse isocourbure (changement de signe \`a $\beta=-1/4$).
\item[{\texttt{[S3]}}] inflation bande viable $f_s=5$--$6$ : $(n_s,r)$ diagonale vs mono (accord) ; $\alpha$, $V^{1/4}$, $\tau_0\approx3{,}5\times10^{-37}$~s, masse du soliton (P1).
\item[{\texttt{[S4]}}] modulation gravitationnelle $\tau_0(R)$ ; $R_c=2\alpha/|\beta|$ ; $\delta\tau_0/\tau_0=-(\beta/4)\Mpl^2R/\alpha$.
\item[{\texttt{[S5]}}] test de la spirale (int\'egration \`a deux champs) : le virage est amorti, $r$ jamais supprim\'e.
\end{itemize}
Fichiers : \texttt{U3\_SYNTH\_reproductible.py} (calculs) et \texttt{Synthese\_figures.py}
(Fig.~\ref{fig:eps}, \ref{fig:pole}, \ref{fig:tau0R}). \emph{Convention} : $\Mpl=1$ ; $\alpha$
s'annule dans les rapports $(n_s,r)$. Le pr\'efixe \texttt{U3\_SYNTH\_} et les codes
\texttt{[S$\cdot$]} sont sp\'ecifiques \`a ce document.

% =====================================================================
\begin{thebibliography}{99}\small
\bibitem{u1} H.\ Bakkaoui, \textit{Champ de dualit\'e d'\'echelle en gravit\'e} (cadre, corrections et r\'esultats), Zenodo (2026).
\bibitem{u2} H.\ Bakkaoui, \textit{Discontinuit\'e d'existence quantique (HDEQ)}, Zenodo (2026).
\bibitem{bardeen} J.\ M.\ Bardeen, Proc.\ GR5 (Tbilisi, 1968), p.\ 174.
\bibitem{hayward} S.\ A.\ Hayward, \textit{Phys.\ Rev.\ Lett.} \textbf{96}, 031103 (2006).
\bibitem{dymnikova} I.\ Dymnikova, \textit{Gen.\ Relativ.\ Gravit.} \textbf{24}, 235 (1992).
\bibitem{penrose} R.\ Penrose, \textit{Phys.\ Rev.\ Lett.} \textbf{14}, 57 (1965).
\bibitem{he} S.\ W.\ Hawking \& G.\ F.\ R.\ Ellis, \textit{The Large Scale Structure of Space-Time} (CUP, 1973).
\bibitem{guth} A.\ H.\ Guth, \textit{Phys.\ Rev.\ D} \textbf{23}, 347 (1981).
\bibitem{ffo} K.\ Freese, J.\ A.\ Frieman \& A.\ V.\ Olinto, \textit{Phys.\ Rev.\ Lett.} \textbf{65}, 3233 (1990).
\bibitem{bek} J.\ D.\ Bekenstein, \textit{Phys.\ Rev.\ D} \textbf{51}, R6608 (1995).
\bibitem{bezr} F.\ L.\ Bezrukov \& M.\ Shaposhnikov, \textit{Phys.\ Lett.\ B} \textbf{659}, 703 (2008).
\bibitem{planck} Planck Collaboration, \textit{Astron.\ Astrophys.} \textbf{641}, A10 (2020).
\bibitem{bk} BICEP/Keck Collaboration, \textit{Phys.\ Rev.\ Lett.} \textbf{127}, 151301 (2021).
\bibitem{litebird} LiteBIRD Collaboration, \textit{Prog.\ Theor.\ Exp.\ Phys.} \textbf{2023}, 042F01 (2023).
\bibitem{cmbs4} K.\ N.\ Abazajian \textit{et al.} (CMB-S4), arXiv:1610.02743 (2016) --- \emph{projet arr\^et\'e, juillet 2025}.
\bibitem{fns} R.\ Z.\ Ferreira, A.\ Notari \& G.\ Simeon, \textit{JCAP} \textbf{11} (2018) 021.
\bibitem{fmm} V.\ P.\ Frolov, M.\ A.\ Markov \& V.\ F.\ Mukhanov, \textit{Phys.\ Rev.\ D} \textbf{41}, 383 (1990).
\bibitem{rv} C.\ Rovelli \& F.\ Vidotto, \textit{Int.\ J.\ Mod.\ Phys.\ D} \textbf{23}, 1442026 (2014).
\bibitem{cr} A.\ Connes \& C.\ Rovelli, \textit{Class.\ Quantum Grav.} \textbf{11}, 2899 (1994).
\bibitem{zurek} W.\ H.\ Zurek, \textit{Rev.\ Mod.\ Phys.} \textbf{75}, 715 (2003).
\bibitem{kaiser} D.\ I.\ Kaiser, \textit{Phys.\ Rev.\ D} \textbf{81}, 084044 (2010).
\bibitem{pw} D.\ N.\ Page \& W.\ K.\ Wootters, \textit{Phys.\ Rev.\ D} \textbf{27}, 2885 (1983).
\bibitem{bars} I.\ Bars, \textit{Class.\ Quantum Grav.} \textbf{18}, 3113 (2001). \statut{\`a v\'erifier}
\bibitem{simeon} G.\ Simeon, arXiv:2002.07625 (2020).
\bibitem{bostan} I.\ D.\ Bostan, arXiv:2209.02434 (2022).
\bibitem{mv} C.\ Marletto \& V.\ Vedral, \textit{Phys.\ Rev.\ D} \textbf{95}, 043510 (2017).
\bibitem{so} Simons Observatory Collaboration (P.\ Ade \textit{et al.}), \textit{JCAP} \textbf{02} (2019) 056, arXiv:1808.07445.
\end{thebibliography}

\end{document}
